% --- Archon physics formalization source begin ---
% archon:physics
% archon:covers PhyXMiniProblems/problem_phyx_mini_0976.lean
% archon:source-report reports/phyx_mini/problem_phyx_mini_0976.source.json

\chapter{Physics problem phyx\_mini\_0976}
\label{ch:PhyXMiniProblems_problem_phyx_mini_0976}

\paragraph{Problem source.}
\#\# Physical scenario

A bar of length \textbackslash{}( L = 0.36\textbackslash{},\textbackslash{}text\{m\} \textbackslash{}) is free to slide without friction on horizontal rails. A uniform magnetic field \textbackslash{}( B = 2.4\textbackslash{},\textbackslash{}text\{T\} \textbackslash{}) is directed into the plane of the figure. At one end of the rails there is a battery with emf \textbackslash{}( \textbackslash{}mathcal\{E\} = 12\textbackslash{},\textbackslash{}text\{V\} \textbackslash{}) and a switch \textbackslash{}( S \textbackslash{}). The bar has mass \textbackslash{}( 0.90\textbackslash{},\textbackslash{}text\{kg\} \textbackslash{}) and resistance \textbackslash{}( 5.0\textbackslash{},\textbackslash{}Omega \textbackslash{}); ignore all other resistance in the circuit. The switch is closed at time \textbackslash{}( t = 0 \textbackslash{}).

\#\# Question

What is the bar's terminal speed?

\#\# Answer choices

- A: A: 12m/s

- B: B: 14m/s

- C: C: 2.4m/s

- D: D: 8.64m/s

\#\# Figure caption (auxiliary; use the image as primary evidence)

The image depicts a physics setup involving a moving rod in a magnetic field.

1. **Objects and Symbols**:
   - **Rod**: A yellow rod is placed horizontally and moves vertically between two parallel rails.
   - **Magnetic Field**: Represented by blue crosses (`X`), indicating the field is directed into the plane of the image.
   - **Battery**: Denoted by `ε` with positive and negative terminals, creating a potential difference.
   - **Switch (`S`)**: Connected in series with the battery.

2. **Connections**:
   - The rod is part of an electrical circuit with connections to the battery and switch.
   - The circuit is completed by the rod moving along the rails.

3. **Text and Symbols**:
   - `\textbackslash{}(\textbackslash{}vec\{B\}\textbackslash{})`: Represents the magnetic field vector going into the page.
   - `\textbackslash{}(L\textbackslash{})`: Indicates the length or distance between the two rails.

4. **Relationships**:
   - The rod is free to move along the rails, influenced by the magnetic field.
   - The battery and switch control the current passing through the circuit when closed.

This setup is typically used to demonstrate electromagnetic induction or Lorentz force concepts in physics.

\#\# Dataset metadata

Electromagnetic Induction; Multi-Formula Reasoning, Predictive Reasoning

\paragraph{Recorded answer/context.}
B: B: 14m/s

\paragraph{Figure/image path.}
/root/proposal\_for\_physic/science-mango/phyx\_mini\_run/phyx\_data/test\_image/976.png

\paragraph{Formalization target.}
create a compiling Lean file with sorry bodies at `PhyXMiniProblems/problem\_phyx\_mini\_0976.lean`. The Lean declarations must preserve the physical quantities, dimensions or dimensional roles, figure labels, governing-law hypotheses, and final relation expressed by this problem.
Use Mathlib/Physlib names found through LeanExplore where available. If a physics API is missing, introduce faithful local abstractions rather than scalar placeholder aliases.

\begin{theorem}[Physics formalization target]
\label{thm:physics:phyx_mini_0976:target}
\lean{PhyXMiniProblems.ProblemPhyXMini0976.problem_phyx_mini_0976}
\uses{def:physics:phyx-mini-0976:phyxminiproblems-problemphyxmini0976-speedinmeterspersecond, def:physics:phyx-mini-0976:phyxminiproblems-problemphyxmini0976-slidingbarrailsetup, def:physics:phyx-mini-0976:phyxminiproblems-problemphyxmini0976-matcheswrittenslidingbarscenario, def:physics:phyx-mini-0976:phyxminiproblems-problemphyxmini0976-matchessuppliedslidingbarfigure, def:physics:phyx-mini-0976:phyxminiproblems-problemphyxmini0976-modelsuniformintopagemagneticfield, def:physics:phyx-mini-0976:phyxminiproblems-problemphyxmini0976-hasslidingbarproblemmeasurements, def:physics:phyx-mini-0976:phyxminiproblems-problemphyxmini0976-hasphysicalslidingbarparameters, def:physics:phyx-mini-0976:phyxminiproblems-problemphyxmini0976-satisfiesbatterydrivenslidingbarlaws, def:physics:phyx-mini-0976:phyxminiproblems-problemphyxmini0976-isterminalspeed, def:physics:phyx-mini-0976:phyxminiproblems-problemphyxmini0976-answerchoice, def:physics:phyx-mini-0976:phyxminiproblems-problemphyxmini0976-isuniqueclosestanswer}
The assigned autoformalize agent should translate this physics problem into Lean declarations in the covered file, with theorem and lemma proof bodies written as `by sorry`.
\end{theorem}
\begin{proof}
This is an autoformalization task, not a proof task. Produce statements that can later be proved by the physics prover without weakening the physical model.
\end{proof}

% --- Archon declaration topology begin ---
\section{Declaration topology}
\begin{definition}[acceleration Dimension]
  \label{def:physics:phyx-mini-0976:phyxminiproblems-problemphyxmini0976-accelerationdimension}
  \lean{PhyXMiniProblems.ProblemPhyXMini0976.accelerationDimension}
  Linear acceleration has physical dimension L T⁻²
\end{definition}
\begin{proof}
  This is the declared model definition of acceleration Dimension.
\end{proof}
\begin{definition}[electric Current Dimension]
  \label{def:physics:phyx-mini-0976:phyxminiproblems-problemphyxmini0976-electriccurrentdimension}
  \lean{PhyXMiniProblems.ProblemPhyXMini0976.electricCurrentDimension}
  Electric current has physical dimension charge per time
\end{definition}
\begin{proof}
  This is the declared model definition of electric Current Dimension.
\end{proof}
\begin{definition}[electromotive Force Dimension]
  \label{def:physics:phyx-mini-0976:phyxminiproblems-problemphyxmini0976-electromotiveforcedimension}
  \lean{PhyXMiniProblems.ProblemPhyXMini0976.electromotiveForceDimension}
  Electromotive force has physical dimension M L² T⁻² C⁻¹
\end{definition}
\begin{proof}
  This is the declared model definition of electromotive Force Dimension.
\end{proof}
\begin{definition}[electrical Resistance Dimension]
  \label{def:physics:phyx-mini-0976:phyxminiproblems-problemphyxmini0976-electricalresistancedimension}
  \lean{PhyXMiniProblems.ProblemPhyXMini0976.electricalResistanceDimension}
  \uses{def:physics:phyx-mini-0976:phyxminiproblems-problemphyxmini0976-electriccurrentdimension, def:physics:phyx-mini-0976:phyxminiproblems-problemphyxmini0976-electromotiveforcedimension}
  Electrical resistance has dimension emf divided by current
\end{definition}
\begin{proof}
  This is the declared model definition of electrical Resistance Dimension.
\end{proof}
\begin{definition}[magnetic Flux Density Dimension]
  \label{def:physics:phyx-mini-0976:phyxminiproblems-problemphyxmini0976-magneticfluxdensitydimension}
  \lean{PhyXMiniProblems.ProblemPhyXMini0976.magneticFluxDensityDimension}
  Magnetic flux density has the tesla dimension M T⁻¹ C⁻¹
\end{definition}
\begin{proof}
  This is the declared model definition of magnetic Flux Density Dimension.
\end{proof}
\begin{definition}[force Dimension]
  \label{def:physics:phyx-mini-0976:phyxminiproblems-problemphyxmini0976-forcedimension}
  \lean{PhyXMiniProblems.ProblemPhyXMini0976.forceDimension}
  Force has the newton dimension M L T⁻²
\end{definition}
\begin{proof}
  This is the declared model definition of force Dimension.
\end{proof}
\begin{definition}[Length Quantity]
  \label{def:physics:phyx-mini-0976:phyxminiproblems-problemphyxmini0976-lengthquantity}
  \lean{PhyXMiniProblems.ProblemPhyXMini0976.LengthQuantity}
  A nonnegative, unit-independent physical length
\end{definition}
\begin{proof}
  This is the declared model definition of Length Quantity.
\end{proof}
\begin{definition}[Duration Quantity]
  \label{def:physics:phyx-mini-0976:phyxminiproblems-problemphyxmini0976-durationquantity}
  \lean{PhyXMiniProblems.ProblemPhyXMini0976.DurationQuantity}
  A nonnegative, unit-independent duration
\end{definition}
\begin{proof}
  This is the declared model definition of Duration Quantity.
\end{proof}
\begin{definition}[Mass Quantity]
  \label{def:physics:phyx-mini-0976:phyxminiproblems-problemphyxmini0976-massquantity}
  \lean{PhyXMiniProblems.ProblemPhyXMini0976.MassQuantity}
  A nonnegative, unit-independent physical mass
\end{definition}
\begin{proof}
  This is the declared model definition of Mass Quantity.
\end{proof}
\begin{definition}[Speed Quantity]
  \label{def:physics:phyx-mini-0976:phyxminiproblems-problemphyxmini0976-speedquantity}
  \lean{PhyXMiniProblems.ProblemPhyXMini0976.SpeedQuantity}
  A nonnegative, unit-independent speed
\end{definition}
\begin{proof}
  This is the declared model definition of Speed Quantity.
\end{proof}
\begin{definition}[Magnetic Flux Density Quantity]
  \label{def:physics:phyx-mini-0976:phyxminiproblems-problemphyxmini0976-magneticfluxdensityquantity}
  \lean{PhyXMiniProblems.ProblemPhyXMini0976.MagneticFluxDensityQuantity}
  \uses{def:physics:phyx-mini-0976:phyxminiproblems-problemphyxmini0976-magneticfluxdensitydimension}
  A nonnegative, unit-independent magnetic-flux-density magnitude
\end{definition}
\begin{proof}
  This is the declared model definition of Magnetic Flux Density Quantity.
\end{proof}
\begin{definition}[Emf Magnitude Quantity]
  \label{def:physics:phyx-mini-0976:phyxminiproblems-problemphyxmini0976-emfmagnitudequantity}
  \lean{PhyXMiniProblems.ProblemPhyXMini0976.EmfMagnitudeQuantity}
  \uses{def:physics:phyx-mini-0976:phyxminiproblems-problemphyxmini0976-electromotiveforcedimension}
  A nonnegative, unit-independent emf magnitude
\end{definition}
\begin{proof}
  This is the declared model definition of Emf Magnitude Quantity.
\end{proof}
\begin{definition}[Resistance Quantity]
  \label{def:physics:phyx-mini-0976:phyxminiproblems-problemphyxmini0976-resistancequantity}
  \lean{PhyXMiniProblems.ProblemPhyXMini0976.ResistanceQuantity}
  \uses{def:physics:phyx-mini-0976:phyxminiproblems-problemphyxmini0976-electricalresistancedimension}
  A nonnegative, unit-independent electrical resistance
\end{definition}
\begin{proof}
  This is the declared model definition of Resistance Quantity.
\end{proof}
\begin{definition}[Signed Current Quantity]
  \label{def:physics:phyx-mini-0976:phyxminiproblems-problemphyxmini0976-signedcurrentquantity}
  \lean{PhyXMiniProblems.ProblemPhyXMini0976.SignedCurrentQuantity}
  \uses{def:physics:phyx-mini-0976:phyxminiproblems-problemphyxmini0976-electriccurrentdimension}
  Signed current, positive in the battery-driven direction through the bar
\end{definition}
\begin{proof}
  This is the declared model definition of Signed Current Quantity.
\end{proof}
\begin{definition}[Signed Force Quantity]
  \label{def:physics:phyx-mini-0976:phyxminiproblems-problemphyxmini0976-signedforcequantity}
  \lean{PhyXMiniProblems.ProblemPhyXMini0976.SignedForceQuantity}
  \uses{def:physics:phyx-mini-0976:phyxminiproblems-problemphyxmini0976-forcedimension}
  Signed horizontal force, positive in the chosen sliding direction
\end{definition}
\begin{proof}
  This is the declared model definition of Signed Force Quantity.
\end{proof}
\begin{definition}[Signed Acceleration Quantity]
  \label{def:physics:phyx-mini-0976:phyxminiproblems-problemphyxmini0976-signedaccelerationquantity}
  \lean{PhyXMiniProblems.ProblemPhyXMini0976.SignedAccelerationQuantity}
  \uses{def:physics:phyx-mini-0976:phyxminiproblems-problemphyxmini0976-accelerationdimension}
  Signed horizontal acceleration, positive in the chosen sliding direction
\end{definition}
\begin{proof}
  This is the declared model definition of Signed Acceleration Quantity.
\end{proof}
\begin{definition}[nonnegative Readout]
  \label{def:physics:phyx-mini-0976:phyxminiproblems-problemphyxmini0976-nonnegativereadout}
  \lean{PhyXMiniProblems.ProblemPhyXMini0976.nonnegativeReadout}
  Read a nonnegative physical quantity in a coherent unit system
\end{definition}
\begin{proof}
  This is the declared model definition of nonnegative Readout.
\end{proof}
\begin{definition}[signed Readout]
  \label{def:physics:phyx-mini-0976:phyxminiproblems-problemphyxmini0976-signedreadout}
  \lean{PhyXMiniProblems.ProblemPhyXMini0976.signedReadout}
  Read a signed physical quantity in a coherent unit system
\end{definition}
\begin{proof}
  This is the declared model definition of signed Readout.
\end{proof}
\begin{definition}[length Readout]
  \label{def:physics:phyx-mini-0976:phyxminiproblems-problemphyxmini0976-lengthreadout}
  \lean{PhyXMiniProblems.ProblemPhyXMini0976.lengthReadout}
  \uses{def:physics:phyx-mini-0976:phyxminiproblems-problemphyxmini0976-lengthquantity, def:physics:phyx-mini-0976:phyxminiproblems-problemphyxmini0976-nonnegativereadout}
  This def records the length Readout component of the problem-specific physical model
\end{definition}
\begin{proof}
  This is the declared model definition of length Readout.
\end{proof}
\begin{definition}[duration Readout]
  \label{def:physics:phyx-mini-0976:phyxminiproblems-problemphyxmini0976-durationreadout}
  \lean{PhyXMiniProblems.ProblemPhyXMini0976.durationReadout}
  \uses{def:physics:phyx-mini-0976:phyxminiproblems-problemphyxmini0976-durationquantity, def:physics:phyx-mini-0976:phyxminiproblems-problemphyxmini0976-nonnegativereadout}
  This def records the duration Readout component of the problem-specific physical model
\end{definition}
\begin{proof}
  This is the declared model definition of duration Readout.
\end{proof}
\begin{definition}[mass Readout]
  \label{def:physics:phyx-mini-0976:phyxminiproblems-problemphyxmini0976-massreadout}
  \lean{PhyXMiniProblems.ProblemPhyXMini0976.massReadout}
  \uses{def:physics:phyx-mini-0976:phyxminiproblems-problemphyxmini0976-massquantity, def:physics:phyx-mini-0976:phyxminiproblems-problemphyxmini0976-nonnegativereadout}
  This def records the mass Readout component of the problem-specific physical model
\end{definition}
\begin{proof}
  This is the declared model definition of mass Readout.
\end{proof}
\begin{definition}[speed Readout]
  \label{def:physics:phyx-mini-0976:phyxminiproblems-problemphyxmini0976-speedreadout}
  \lean{PhyXMiniProblems.ProblemPhyXMini0976.speedReadout}
  \uses{def:physics:phyx-mini-0976:phyxminiproblems-problemphyxmini0976-speedquantity, def:physics:phyx-mini-0976:phyxminiproblems-problemphyxmini0976-nonnegativereadout}
  This def records the speed Readout component of the problem-specific physical model
\end{definition}
\begin{proof}
  This is the declared model definition of speed Readout.
\end{proof}
\begin{definition}[magnetic Flux Density Readout]
  \label{def:physics:phyx-mini-0976:phyxminiproblems-problemphyxmini0976-magneticfluxdensityreadout}
  \lean{PhyXMiniProblems.ProblemPhyXMini0976.magneticFluxDensityReadout}
  \uses{def:physics:phyx-mini-0976:phyxminiproblems-problemphyxmini0976-magneticfluxdensityquantity, def:physics:phyx-mini-0976:phyxminiproblems-problemphyxmini0976-nonnegativereadout}
  This def records the magnetic Flux Density Readout component of the problem-specific physical model
\end{definition}
\begin{proof}
  This is the declared model definition of magnetic Flux Density Readout.
\end{proof}
\begin{definition}[emf Readout]
  \label{def:physics:phyx-mini-0976:phyxminiproblems-problemphyxmini0976-emfreadout}
  \lean{PhyXMiniProblems.ProblemPhyXMini0976.emfReadout}
  \uses{def:physics:phyx-mini-0976:phyxminiproblems-problemphyxmini0976-emfmagnitudequantity, def:physics:phyx-mini-0976:phyxminiproblems-problemphyxmini0976-nonnegativereadout}
  This def records the emf Readout component of the problem-specific physical model
\end{definition}
\begin{proof}
  This is the declared model definition of emf Readout.
\end{proof}
\begin{definition}[resistance Readout]
  \label{def:physics:phyx-mini-0976:phyxminiproblems-problemphyxmini0976-resistancereadout}
  \lean{PhyXMiniProblems.ProblemPhyXMini0976.resistanceReadout}
  \uses{def:physics:phyx-mini-0976:phyxminiproblems-problemphyxmini0976-resistancequantity, def:physics:phyx-mini-0976:phyxminiproblems-problemphyxmini0976-nonnegativereadout}
  This def records the resistance Readout component of the problem-specific physical model
\end{definition}
\begin{proof}
  This is the declared model definition of resistance Readout.
\end{proof}
\begin{definition}[current Readout]
  \label{def:physics:phyx-mini-0976:phyxminiproblems-problemphyxmini0976-currentreadout}
  \lean{PhyXMiniProblems.ProblemPhyXMini0976.currentReadout}
  \uses{def:physics:phyx-mini-0976:phyxminiproblems-problemphyxmini0976-signedcurrentquantity, def:physics:phyx-mini-0976:phyxminiproblems-problemphyxmini0976-signedreadout}
  This def records the current Readout component of the problem-specific physical model
\end{definition}
\begin{proof}
  This is the declared model definition of current Readout.
\end{proof}
\begin{definition}[force Readout]
  \label{def:physics:phyx-mini-0976:phyxminiproblems-problemphyxmini0976-forcereadout}
  \lean{PhyXMiniProblems.ProblemPhyXMini0976.forceReadout}
  \uses{def:physics:phyx-mini-0976:phyxminiproblems-problemphyxmini0976-signedforcequantity, def:physics:phyx-mini-0976:phyxminiproblems-problemphyxmini0976-signedreadout}
  This def records the force Readout component of the problem-specific physical model
\end{definition}
\begin{proof}
  This is the declared model definition of force Readout.
\end{proof}
\begin{definition}[acceleration Readout]
  \label{def:physics:phyx-mini-0976:phyxminiproblems-problemphyxmini0976-accelerationreadout}
  \lean{PhyXMiniProblems.ProblemPhyXMini0976.accelerationReadout}
  \uses{def:physics:phyx-mini-0976:phyxminiproblems-problemphyxmini0976-signedaccelerationquantity, def:physics:phyx-mini-0976:phyxminiproblems-problemphyxmini0976-signedreadout}
  This def records the acceleration Readout component of the problem-specific physical model
\end{definition}
\begin{proof}
  This is the declared model definition of acceleration Readout.
\end{proof}
\begin{definition}[length In Meters]
  \label{def:physics:phyx-mini-0976:phyxminiproblems-problemphyxmini0976-lengthinmeters}
  \lean{PhyXMiniProblems.ProblemPhyXMini0976.lengthInMeters}
  \uses{def:physics:phyx-mini-0976:phyxminiproblems-problemphyxmini0976-lengthquantity, def:physics:phyx-mini-0976:phyxminiproblems-problemphyxmini0976-lengthreadout}
  SI readout of length, in metres
\end{definition}
\begin{proof}
  This is the declared model definition of length In Meters.
\end{proof}
\begin{definition}[duration In Seconds]
  \label{def:physics:phyx-mini-0976:phyxminiproblems-problemphyxmini0976-durationinseconds}
  \lean{PhyXMiniProblems.ProblemPhyXMini0976.durationInSeconds}
  \uses{def:physics:phyx-mini-0976:phyxminiproblems-problemphyxmini0976-durationquantity, def:physics:phyx-mini-0976:phyxminiproblems-problemphyxmini0976-durationreadout}
  SI readout of duration, in seconds
\end{definition}
\begin{proof}
  This is the declared model definition of duration In Seconds.
\end{proof}
\begin{definition}[mass In Kilograms]
  \label{def:physics:phyx-mini-0976:phyxminiproblems-problemphyxmini0976-massinkilograms}
  \lean{PhyXMiniProblems.ProblemPhyXMini0976.massInKilograms}
  \uses{def:physics:phyx-mini-0976:phyxminiproblems-problemphyxmini0976-massquantity, def:physics:phyx-mini-0976:phyxminiproblems-problemphyxmini0976-massreadout}
  SI readout of mass, in kilograms
\end{definition}
\begin{proof}
  This is the declared model definition of mass In Kilograms.
\end{proof}
\begin{definition}[speed In Meters Per Second]
  \label{def:physics:phyx-mini-0976:phyxminiproblems-problemphyxmini0976-speedinmeterspersecond}
  \lean{PhyXMiniProblems.ProblemPhyXMini0976.speedInMetersPerSecond}
  \uses{def:physics:phyx-mini-0976:phyxminiproblems-problemphyxmini0976-speedquantity, def:physics:phyx-mini-0976:phyxminiproblems-problemphyxmini0976-speedreadout}
  SI readout of speed, in metres per second
\end{definition}
\begin{proof}
  This is the declared model definition of speed In Meters Per Second.
\end{proof}
\begin{definition}[magnetic Flux Density In Teslas]
  \label{def:physics:phyx-mini-0976:phyxminiproblems-problemphyxmini0976-magneticfluxdensityinteslas}
  \lean{PhyXMiniProblems.ProblemPhyXMini0976.magneticFluxDensityInTeslas}
  \uses{def:physics:phyx-mini-0976:phyxminiproblems-problemphyxmini0976-magneticfluxdensityquantity, def:physics:phyx-mini-0976:phyxminiproblems-problemphyxmini0976-magneticfluxdensityreadout}
  SI readout of magnetic flux density, in teslas
\end{definition}
\begin{proof}
  This is the declared model definition of magnetic Flux Density In Teslas.
\end{proof}
\begin{definition}[emf In Volts]
  \label{def:physics:phyx-mini-0976:phyxminiproblems-problemphyxmini0976-emfinvolts}
  \lean{PhyXMiniProblems.ProblemPhyXMini0976.emfInVolts}
  \uses{def:physics:phyx-mini-0976:phyxminiproblems-problemphyxmini0976-emfmagnitudequantity, def:physics:phyx-mini-0976:phyxminiproblems-problemphyxmini0976-emfreadout}
  SI readout of emf, in volts
\end{definition}
\begin{proof}
  This is the declared model definition of emf In Volts.
\end{proof}
\begin{definition}[resistance In Ohms]
  \label{def:physics:phyx-mini-0976:phyxminiproblems-problemphyxmini0976-resistanceinohms}
  \lean{PhyXMiniProblems.ProblemPhyXMini0976.resistanceInOhms}
  \uses{def:physics:phyx-mini-0976:phyxminiproblems-problemphyxmini0976-resistancequantity, def:physics:phyx-mini-0976:phyxminiproblems-problemphyxmini0976-resistancereadout}
  SI readout of resistance, in ohms
\end{definition}
\begin{proof}
  This is the declared model definition of resistance In Ohms.
\end{proof}
\begin{definition}[current In Amperes]
  \label{def:physics:phyx-mini-0976:phyxminiproblems-problemphyxmini0976-currentinamperes}
  \lean{PhyXMiniProblems.ProblemPhyXMini0976.currentInAmperes}
  \uses{def:physics:phyx-mini-0976:phyxminiproblems-problemphyxmini0976-signedcurrentquantity, def:physics:phyx-mini-0976:phyxminiproblems-problemphyxmini0976-currentreadout}
  SI readout of signed current, in amperes
\end{definition}
\begin{proof}
  This is the declared model definition of current In Amperes.
\end{proof}
\begin{definition}[force In Newtons]
  \label{def:physics:phyx-mini-0976:phyxminiproblems-problemphyxmini0976-forceinnewtons}
  \lean{PhyXMiniProblems.ProblemPhyXMini0976.forceInNewtons}
  \uses{def:physics:phyx-mini-0976:phyxminiproblems-problemphyxmini0976-signedforcequantity, def:physics:phyx-mini-0976:phyxminiproblems-problemphyxmini0976-forcereadout}
  SI readout of signed force, in newtons
\end{definition}
\begin{proof}
  This is the declared model definition of force In Newtons.
\end{proof}
\begin{definition}[acceleration In Meters Per Second Squared]
  \label{def:physics:phyx-mini-0976:phyxminiproblems-problemphyxmini0976-accelerationinmeterspersecondsquared}
  \lean{PhyXMiniProblems.ProblemPhyXMini0976.accelerationInMetersPerSecondSquared}
  \uses{def:physics:phyx-mini-0976:phyxminiproblems-problemphyxmini0976-signedaccelerationquantity, def:physics:phyx-mini-0976:phyxminiproblems-problemphyxmini0976-accelerationreadout}
  SI readout of signed acceleration, in metres per second squared
\end{definition}
\begin{proof}
  This is the declared model definition of acceleration In Meters Per Second Squared.
\end{proof}
\begin{definition}[Apparatus Component]
  \label{def:physics:phyx-mini-0976:phyxminiproblems-problemphyxmini0976-apparatuscomponent}
  \lean{PhyXMiniProblems.ProblemPhyXMini0976.ApparatusComponent}
  Components visibly present in the supplied apparatus figure
\end{definition}
\begin{proof}
  This is the declared model definition of Apparatus Component.
\end{proof}
\begin{definition}[Figure Label]
  \label{def:physics:phyx-mini-0976:phyxminiproblems-problemphyxmini0976-figurelabel}
  \lean{PhyXMiniProblems.ProblemPhyXMini0976.FigureLabel}
  Symbolic labels printed in the supplied figure
\end{definition}
\begin{proof}
  This is the declared model definition of Figure Label.
\end{proof}
\begin{definition}[Figure Direction]
  \label{def:physics:phyx-mini-0976:phyxminiproblems-problemphyxmini0976-figuredirection}
  \lean{PhyXMiniProblems.ProblemPhyXMini0976.FigureDirection}
  Directions distinguished by the drawing and its page normal
\end{definition}
\begin{proof}
  This is the declared model definition of Figure Direction.
\end{proof}
\begin{definition}[Figure Axis]
  \label{def:physics:phyx-mini-0976:phyxminiproblems-problemphyxmini0976-figureaxis}
  \lean{PhyXMiniProblems.ProblemPhyXMini0976.FigureAxis}
  This inductive records the Figure Axis component of the problem-specific physical model
\end{definition}
\begin{proof}
  This is the declared model definition of Figure Axis.
\end{proof}
\begin{definition}[Magnetic Field Glyph]
  \label{def:physics:phyx-mini-0976:phyxminiproblems-problemphyxmini0976-magneticfieldglyph}
  \lean{PhyXMiniProblems.ProblemPhyXMini0976.MagneticFieldGlyph}
  Crosses in the raster encode a field directed into the page
\end{definition}
\begin{proof}
  This is the declared model definition of Magnetic Field Glyph.
\end{proof}
\begin{definition}[Switch State]
  \label{def:physics:phyx-mini-0976:phyxminiproblems-problemphyxmini0976-switchstate}
  \lean{PhyXMiniProblems.ProblemPhyXMini0976.SwitchState}
  This inductive records the Switch State component of the problem-specific physical model
\end{definition}
\begin{proof}
  This is the declared model definition of Switch State.
\end{proof}
\begin{definition}[Rail Contact Model]
  \label{def:physics:phyx-mini-0976:phyxminiproblems-problemphyxmini0976-railcontactmodel}
  \lean{PhyXMiniProblems.ProblemPhyXMini0976.RailContactModel}
  This inductive records the Rail Contact Model component of the problem-specific physical model
\end{definition}
\begin{proof}
  This is the declared model definition of Rail Contact Model.
\end{proof}
\begin{definition}[Sliding Bar Geometry]
  \label{def:physics:phyx-mini-0976:phyxminiproblems-problemphyxmini0976-slidingbargeometry}
  \lean{PhyXMiniProblems.ProblemPhyXMini0976.SlidingBarGeometry}
  This inductive records the Sliding Bar Geometry component of the problem-specific physical model
\end{definition}
\begin{proof}
  This is the declared model definition of Sliding Bar Geometry.
\end{proof}
\begin{definition}[Circuit Resistance Model]
  \label{def:physics:phyx-mini-0976:phyxminiproblems-problemphyxmini0976-circuitresistancemodel}
  \lean{PhyXMiniProblems.ProblemPhyXMini0976.CircuitResistanceModel}
  This inductive records the Circuit Resistance Model component of the problem-specific physical model
\end{definition}
\begin{proof}
  This is the declared model definition of Circuit Resistance Model.
\end{proof}
\begin{definition}[Magnetic Field Uniformity]
  \label{def:physics:phyx-mini-0976:phyxminiproblems-problemphyxmini0976-magneticfielduniformity}
  \lean{PhyXMiniProblems.ProblemPhyXMini0976.MagneticFieldUniformity}
  This inductive records the Magnetic Field Uniformity component of the problem-specific physical model
\end{definition}
\begin{proof}
  This is the declared model definition of Magnetic Field Uniformity.
\end{proof}
\begin{definition}[Sliding Bar Rail Figure]
  \label{def:physics:phyx-mini-0976:phyxminiproblems-problemphyxmini0976-slidingbarrailfigure}
  \lean{PhyXMiniProblems.ProblemPhyXMini0976.SlidingBarRailFigure}
  \uses{def:physics:phyx-mini-0976:phyxminiproblems-problemphyxmini0976-apparatuscomponent, def:physics:phyx-mini-0976:phyxminiproblems-problemphyxmini0976-figurelabel, def:physics:phyx-mini-0976:phyxminiproblems-problemphyxmini0976-figuredirection, def:physics:phyx-mini-0976:phyxminiproblems-problemphyxmini0976-figureaxis, def:physics:phyx-mini-0976:phyxminiproblems-problemphyxmini0976-magneticfieldglyph, def:physics:phyx-mini-0976:phyxminiproblems-problemphyxmini0976-switchstate}
  Literal presentation data in 976.png. The switch is drawn open, matching the instant just before its stated closure at t = 0. The raster contains no motion arrow and no terminal-speed readout
\end{definition}
\begin{proof}
  This is the declared model definition of Sliding Bar Rail Figure.
\end{proof}
\begin{definition}[Direction Vector]
  \label{def:physics:phyx-mini-0976:phyxminiproblems-problemphyxmini0976-directionvector}
  \lean{PhyXMiniProblems.ProblemPhyXMini0976.DirectionVector}
  A dimensionless direction vector in physical three-space
\end{definition}
\begin{proof}
  This is the declared model definition of Direction Vector.
\end{proof}
\begin{definition}[into Page Unit Vector]
  \label{def:physics:phyx-mini-0976:phyxminiproblems-problemphyxmini0976-intopageunitvector}
  \lean{PhyXMiniProblems.ProblemPhyXMini0976.intoPageUnitVector}
  \uses{def:physics:phyx-mini-0976:phyxminiproblems-problemphyxmini0976-directionvector}
  Unit vector pointing into the page
\end{definition}
\begin{proof}
  This is the declared model definition of into Page Unit Vector.
\end{proof}
\begin{definition}[Sliding Bar Rail Setup]
  \label{def:physics:phyx-mini-0976:phyxminiproblems-problemphyxmini0976-slidingbarrailsetup}
  \lean{PhyXMiniProblems.ProblemPhyXMini0976.SlidingBarRailSetup}
  \uses{def:physics:phyx-mini-0976:phyxminiproblems-problemphyxmini0976-lengthquantity, def:physics:phyx-mini-0976:phyxminiproblems-problemphyxmini0976-durationquantity, def:physics:phyx-mini-0976:phyxminiproblems-problemphyxmini0976-massquantity, def:physics:phyx-mini-0976:phyxminiproblems-problemphyxmini0976-speedquantity, def:physics:phyx-mini-0976:phyxminiproblems-problemphyxmini0976-magneticfluxdensityquantity, def:physics:phyx-mini-0976:phyxminiproblems-problemphyxmini0976-emfmagnitudequantity, def:physics:phyx-mini-0976:phyxminiproblems-problemphyxmini0976-resistancequantity, def:physics:phyx-mini-0976:phyxminiproblems-problemphyxmini0976-signedcurrentquantity, def:physics:phyx-mini-0976:phyxminiproblems-problemphyxmini0976-signedforcequantity, def:physics:phyx-mini-0976:phyxminiproblems-problemphyxmini0976-signedaccelerationquantity, def:physics:phyx-mini-0976:phyxminiproblems-problemphyxmini0976-figuredirection, def:physics:phyx-mini-0976:phyxminiproblems-problemphyxmini0976-switchstate, def:physics:phyx-mini-0976:phyxminiproblems-problemphyxmini0976-railcontactmodel, def:physics:phyx-mini-0976:phyxminiproblems-problemphyxmini0976-slidingbargeometry, def:physics:phyx-mini-0976:phyxminiproblems-problemphyxmini0976-circuitresistancemodel, def:physics:phyx-mini-0976:phyxminiproblems-problemphyxmini0976-magneticfielduniformity, def:physics:phyx-mini-0976:phyxminiproblems-problemphyxmini0976-slidingbarrailfigure}
  Independent apparatus data and speed-dependent observables. In particular, terminalSpeed is a physical speed rather than a scalar answer value, and the functions for emf, current, force, and acceleration are not defined from the requested conclusion
\end{definition}
\begin{proof}
  This is the declared model definition of Sliding Bar Rail Setup.
\end{proof}
\begin{definition}[Matches Written Sliding Bar Scenario]
  \label{def:physics:phyx-mini-0976:phyxminiproblems-problemphyxmini0976-matcheswrittenslidingbarscenario}
  \lean{PhyXMiniProblems.ProblemPhyXMini0976.MatchesWrittenSlidingBarScenario}
  \uses{def:physics:phyx-mini-0976:phyxminiproblems-problemphyxmini0976-slidingbarrailsetup}
  Qualitative apparatus facts stated in the problem prose
\end{definition}
\begin{proof}
  This is the declared model definition of Matches Written Sliding Bar Scenario.
\end{proof}
\begin{definition}[Matches Supplied Sliding Bar Figure]
  \label{def:physics:phyx-mini-0976:phyxminiproblems-problemphyxmini0976-matchessuppliedslidingbarfigure}
  \lean{PhyXMiniProblems.ProblemPhyXMini0976.MatchesSuppliedSlidingBarFigure}
  \uses{def:physics:phyx-mini-0976:phyxminiproblems-problemphyxmini0976-slidingbarrailsetup}
  Literal geometry, labels, polarity, field glyphs, and switch state in the image
\end{definition}
\begin{proof}
  This is the declared model definition of Matches Supplied Sliding Bar Figure.
\end{proof}
\begin{definition}[Models Uniform Into Page Magnetic Field]
  \label{def:physics:phyx-mini-0976:phyxminiproblems-problemphyxmini0976-modelsuniformintopagemagneticfield}
  \lean{PhyXMiniProblems.ProblemPhyXMini0976.ModelsUniformIntoPageMagneticField}
  \uses{def:physics:phyx-mini-0976:phyxminiproblems-problemphyxmini0976-magneticfluxdensityinteslas, def:physics:phyx-mini-0976:phyxminiproblems-problemphyxmini0976-intopageunitvector, def:physics:phyx-mini-0976:phyxminiproblems-problemphyxmini0976-slidingbarrailsetup}
  The Physlib spacetime-dependent field is constant, points into the page, and has norm calibrated by the independent dimensionful tesla magnitude
\end{definition}
\begin{proof}
  This is the declared model definition of Models Uniform Into Page Magnetic Field.
\end{proof}
\begin{definition}[Has Sliding Bar Problem Measurements]
  \label{def:physics:phyx-mini-0976:phyxminiproblems-problemphyxmini0976-hasslidingbarproblemmeasurements}
  \lean{PhyXMiniProblems.ProblemPhyXMini0976.HasSlidingBarProblemMeasurements}
  \uses{def:physics:phyx-mini-0976:phyxminiproblems-problemphyxmini0976-lengthinmeters, def:physics:phyx-mini-0976:phyxminiproblems-problemphyxmini0976-durationinseconds, def:physics:phyx-mini-0976:phyxminiproblems-problemphyxmini0976-massinkilograms, def:physics:phyx-mini-0976:phyxminiproblems-problemphyxmini0976-magneticfluxdensityinteslas, def:physics:phyx-mini-0976:phyxminiproblems-problemphyxmini0976-emfinvolts, def:physics:phyx-mini-0976:phyxminiproblems-problemphyxmini0976-resistanceinohms, def:physics:phyx-mini-0976:phyxminiproblems-problemphyxmini0976-slidingbarrailsetup}
  Numerical readouts stated by the problem. Resistance outside the bar is zero because it is explicitly ignored. No speed, current, force, or acceleration readout occurs here
\end{definition}
\begin{proof}
  This is the declared model definition of Has Sliding Bar Problem Measurements.
\end{proof}
\begin{definition}[Has Physical Sliding Bar Parameters]
  \label{def:physics:phyx-mini-0976:phyxminiproblems-problemphyxmini0976-hasphysicalslidingbarparameters}
  \lean{PhyXMiniProblems.ProblemPhyXMini0976.HasPhysicalSlidingBarParameters}
  \uses{def:physics:phyx-mini-0976:phyxminiproblems-problemphyxmini0976-lengthinmeters, def:physics:phyx-mini-0976:phyxminiproblems-problemphyxmini0976-massinkilograms, def:physics:phyx-mini-0976:phyxminiproblems-problemphyxmini0976-magneticfluxdensityinteslas, def:physics:phyx-mini-0976:phyxminiproblems-problemphyxmini0976-emfinvolts, def:physics:phyx-mini-0976:phyxminiproblems-problemphyxmini0976-resistanceinohms, def:physics:phyx-mini-0976:phyxminiproblems-problemphyxmini0976-slidingbarrailsetup}
  Positivity and nondegeneracy conditions needed by the physical laws
\end{definition}
\begin{proof}
  This is the declared model definition of Has Physical Sliding Bar Parameters.
\end{proof}
\begin{definition}[Satisfies Battery Driven Sliding Bar Laws]
  \label{def:physics:phyx-mini-0976:phyxminiproblems-problemphyxmini0976-satisfiesbatterydrivenslidingbarlaws}
  \lean{PhyXMiniProblems.ProblemPhyXMini0976.SatisfiesBatteryDrivenSlidingBarLaws}
  \uses{def:physics:phyx-mini-0976:phyxminiproblems-problemphyxmini0976-speedquantity, def:physics:phyx-mini-0976:phyxminiproblems-problemphyxmini0976-lengthreadout, def:physics:phyx-mini-0976:phyxminiproblems-problemphyxmini0976-massreadout, def:physics:phyx-mini-0976:phyxminiproblems-problemphyxmini0976-speedreadout, def:physics:phyx-mini-0976:phyxminiproblems-problemphyxmini0976-magneticfluxdensityreadout, def:physics:phyx-mini-0976:phyxminiproblems-problemphyxmini0976-emfreadout, def:physics:phyx-mini-0976:phyxminiproblems-problemphyxmini0976-resistancereadout, def:physics:phyx-mini-0976:phyxminiproblems-problemphyxmini0976-currentreadout, def:physics:phyx-mini-0976:phyxminiproblems-problemphyxmini0976-forcereadout, def:physics:phyx-mini-0976:phyxminiproblems-problemphyxmini0976-accelerationreadout, def:physics:phyx-mini-0976:phyxminiproblems-problemphyxmini0976-slidingbarrailsetup}
  Governing relations for every candidate speed. They do not assign the terminal speed or assert the terminal emf balance as a premise
\end{definition}
\begin{proof}
  This is the declared model definition of Satisfies Battery Driven Sliding Bar Laws.
\end{proof}
\begin{definition}[Is Terminal Speed]
  \label{def:physics:phyx-mini-0976:phyxminiproblems-problemphyxmini0976-isterminalspeed}
  \lean{PhyXMiniProblems.ProblemPhyXMini0976.IsTerminalSpeed}
  \uses{def:physics:phyx-mini-0976:phyxminiproblems-problemphyxmini0976-speedquantity, def:physics:phyx-mini-0976:phyxminiproblems-problemphyxmini0976-accelerationinmeterspersecondsquared, def:physics:phyx-mini-0976:phyxminiproblems-problemphyxmini0976-slidingbarrailsetup}
  A terminal speed is characterized dynamically by zero acceleration. This definition contains no numerical speed and does not unfold to the answer
\end{definition}
\begin{proof}
  This is the declared model definition of Is Terminal Speed.
\end{proof}
\begin{lemma}[terminal Speed eq battery Emf div field length]
  \label{lem:physics:phyx-mini-0976:phyxminiproblems-problemphyxmini0976-terminalspeed-eq-batteryemf-div-field-length}
  \lean{PhyXMiniProblems.ProblemPhyXMini0976.terminalSpeed_eq_batteryEmf_div_field_length}
  \uses{def:physics:phyx-mini-0976:phyxminiproblems-problemphyxmini0976-lengthinmeters, def:physics:phyx-mini-0976:phyxminiproblems-problemphyxmini0976-speedinmeterspersecond, def:physics:phyx-mini-0976:phyxminiproblems-problemphyxmini0976-magneticfluxdensityinteslas, def:physics:phyx-mini-0976:phyxminiproblems-problemphyxmini0976-emfinvolts, def:physics:phyx-mini-0976:phyxminiproblems-problemphyxmini0976-slidingbarrailsetup, def:physics:phyx-mini-0976:phyxminiproblems-problemphyxmini0976-hasphysicalslidingbarparameters, def:physics:phyx-mini-0976:phyxminiproblems-problemphyxmini0976-satisfiesbatterydrivenslidingbarlaws, def:physics:phyx-mini-0976:phyxminiproblems-problemphyxmini0976-isterminalspeed}
  At zero acceleration, force and current vanish. The loop law then balances the battery and motional emfs, yielding the general terminal-speed formula
\end{lemma}
\begin{proof}
  The conclusion follows from the governing relations and figure readouts named in the statement of terminal Speed eq battery Emf div field length.
\end{proof}
\begin{definition}[Answer Choice]
  \label{def:physics:phyx-mini-0976:phyxminiproblems-problemphyxmini0976-answerchoice}
  \lean{PhyXMiniProblems.ProblemPhyXMini0976.AnswerChoice}
  Labels of the four terminal-speed choices printed with the problem
\end{definition}
\begin{proof}
  This is the declared model definition of Answer Choice.
\end{proof}
\begin{definition}[displayed Speed In Meters Per Second]
  \label{def:physics:phyx-mini-0976:phyxminiproblems-problemphyxmini0976-answerchoice-displayedspeedinmeterspersecond}
  \lean{PhyXMiniProblems.ProblemPhyXMini0976.AnswerChoice.displayedSpeedInMetersPerSecond}
  \uses{def:physics:phyx-mini-0976:phyxminiproblems-problemphyxmini0976-answerchoice}
  Speed in metres per second printed beside each answer label
\end{definition}
\begin{proof}
  This is the declared model definition of displayed Speed In Meters Per Second.
\end{proof}
\begin{definition}[recorded Dataset Answer]
  \label{def:physics:phyx-mini-0976:phyxminiproblems-problemphyxmini0976-recordeddatasetanswer}
  \lean{PhyXMiniProblems.ProblemPhyXMini0976.recordedDatasetAnswer}
  \uses{def:physics:phyx-mini-0976:phyxminiproblems-problemphyxmini0976-answerchoice}
  Dataset answer metadata, retained separately and never used as a premise
\end{definition}
\begin{proof}
  This is the declared model definition of recorded Dataset Answer.
\end{proof}
\begin{definition}[Is Unique Closest Answer]
  \label{def:physics:phyx-mini-0976:phyxminiproblems-problemphyxmini0976-isuniqueclosestanswer}
  \lean{PhyXMiniProblems.ProblemPhyXMini0976.IsUniqueClosestAnswer}
  \uses{def:physics:phyx-mini-0976:phyxminiproblems-problemphyxmini0976-speedinmeterspersecond, def:physics:phyx-mini-0976:phyxminiproblems-problemphyxmini0976-slidingbarrailsetup, def:physics:phyx-mini-0976:phyxminiproblems-problemphyxmini0976-answerchoice}
  A displayed answer is strictly closer to the physical speed than every rival
\end{definition}
\begin{proof}
  This is the declared model definition of Is Unique Closest Answer.
\end{proof}
% --- Archon declaration topology end ---
% --- Archon physics formalization source end ---
